% Optional precise replacement for the overlarge-domain reduction in H9(b).
% Deduction in run JAC1964-20261121. /root/h9_height.
% This is conditional exactly where stated. Root whole-read the final proof,
% including the accepted032 upgrade, and reported PASS on 2026-11-21.

\subsection*{H9-T11. A conditional reduction using an elliptic product}

The invalid map from all homologically trivial hyperplane restrictions
to an incidence Jacobian is unnecessary for the following more limited
reduction. It identifies separately the positivity input and the remaining
numerical-versus-homological gate.

\begin{proposition}[H9-T11]
Let $X/k$ be smooth projective of dimension $2m$ over an algebraically
closed field, let $h$ be an ample class, and let $E/k$ be an elliptic
curve with origin $0$. Suppose that the primitive projector for the
middle odd cohomology of $V=X\times E$, for the ample class
$h+[0]_E$, is algebraic. Suppose the primitive incidence assertion of
Theorem 4 holds for $V$: its signed incidence class is a rational
polarization on $P^{m+1}(V)_{\rm prim}$.

Then the form
\[
 q_X(Z)=(-1)^m\deg_X(Z^2)
\]
is nonnegative on all rational algebraic cycles $Z$ with
$h\,\operatorname{cl}(Z)=0$. On this primitive domain its zero space
is exactly the subgroup of numerically trivial cycles. Consequently
it is positive definite on the quotient of that domain by numerical
equivalence. If numerical and homological equivalence coincide for
middle cycles on $X$, it is positive definite on the primitive
algebraic cohomology classes: this is the middle Hodge-standard
assertion for $(X,h)$.
\end{proposition}

\begin{proof}
On $E\times E$ use the normalized Poincar\'e divisor
\[
 P_E=\Delta_E-E\times\{0\}-\{0\}\times E.
\]
As a rational Chow correspondence it is the degree-one projector:
$P_E^t\circ P_E=P_E$. Its diagonal pullback has degree $-2$,
because the diagonal normal bundle is $T_E$ of degree zero and the
two axis terms each restrict to $[0]$. Thus
\[
 \tfrac12\Delta_E^*P_E=-[0]
       \quad\text{in }\operatorname{NS}(E)_{\mathbf Q}.       \tag{H9.36}
\]

For a codimension-$m$ cycle $Z$ on $X$, form the pointed family
\[
 z_Z=Z\boxtimes P_E
    \in\operatorname{CH}^{m+1}(E\times(X\times E))_{\mathbf Q},
\]
with the first copy of $E$ the parameter. Every fibre difference is
algebraically trivial. Its relevant class is in
$H^1(E)\otimes H^{2m}(X)\otimes H^1(E)$.
If $h\operatorname{cl}(Z)=0$, then this class is primitive on $V$:
the $h$ term kills its $X$ component, and $[0]_E H^1(E)=0$.
The algebraic primitive projector and incidence detection (JAC-P01)
therefore show that its family homomorphism factors through
\[
 a_Z:E\longrightarrow B:=P^{m+1}(V)_{\rm prim}
\]
in the rational category. Clear positive denominators whenever an
actual homomorphism is needed.

Direct composition of the two family correspondences gives the
self-incidence divisor
\[
 D_Z=\deg_X(Z^2)P_E.
\]
If $\ell_B$ is the normalized signed incidence class on $B$, the
incidence identity and (H9.36) give
\[
 a_Z^*\ell_B
   =\frac{(-1)^{m+1}}2\Delta_E^*D_Z
   =(-1)^m\deg_X(Z^2)[0].                                  \tag{H9.37}
\]
The pullback of a rational polarization is nef, and its degree on
the elliptic curve is zero if and only if $a_Z=0$. Indeed a nonzero
homomorphism from an elliptic curve has finite kernel and its image
is an elliptic subvariety on which an ample line bundle has positive
degree. Thus $q_X(Z)\ge0$, and equality is equivalent to $a_Z=0$.

If $a_Z=0$, test the family against $z_{Z'}$ for an arbitrary rational
codimension-$m$ cycle $Z'$ on $X$. Incidence equivalence says that
the resulting divisor correspondence on the two parameter elliptic
curves induces the zero homomorphism. That correspondence is
\[
 \deg_X(ZZ')P_E.
\]
Since $P_E$ induces a nonzero isomorphism $E\to\operatorname{Pic}^0(E)$,
one has $\deg_X(ZZ')=0$ for every $Z'$. This is precisely numerical
triviality of $Z$. Conversely a numerically trivial $Z$ has $Z^2=0$.
We have proved the exact zero-space assertion. Positivity of the
induced rational bilinear form on the numerical quotient follows by
polarization. If middle numerical and homological equivalence coincide,
this quotient is exactly the space of primitive algebraic cohomology
classes, which proves the final assertion.
\end{proof}

\paragraph{Exact inputs and the lower-degree consequence.}
The algebraic-projector hypothesis holds, for example, if the
Lefschetz standard conjecture holds for $X$, since it then holds for
$X\times E$. The middle equality of numerical and homological
equivalence is a separate input; it has not been inferred merely from
positivity on an incidence quotient.

Here is a precise way to recover another degree from middle degree.
Let $X_0$ have dimension $n$, let $1\le p\le n/2$, and choose a smooth
complete intersection $i:W\hookrightarrow X_0$ of $n-2p$ divisors
in a sufficiently ample multiple of a chosen class $h_0$, with
$\dim W=2p$. Write their common class as $h=d h_0$.
For a primitive class $z\in\operatorname{Alg}^p(X_0)$, meaning
$h^{n-2p+1}z=0$, its restriction $i^*z$ is primitive on $W$.
Indeed $i_*L_h i^*z=h^{n-2p+1}z=0$, and
\[
 i_*:H^{2p+2}(W)\longrightarrow H^{2n-2p+2}(X_0)
\]
is injective by dual weak Lefschetz. Weak Lefschetz also makes
$i^*:H^{2p}(X_0)\to H^{2p}(W)$ injective. Finally
\[
 \deg_W((i^*z)^2)=\deg_{X_0}(z^2h^{n-2p}).
\]
Thus the middle assertion for $W$ implies the degree-$p$ assertion
for $X_0$. H9-T11 supplies that middle assertion from primitive
Theorem 4 for the odd variety $W\times E$, together with its displayed
algebraic-projector hypothesis and the middle numerical/homological
equality for $W$. These are the precise hypotheses for this reduction.
They must not be replaced by the assertion that a merely homologically
trivial hyperplane restriction defines a point of $P^p$.

\paragraph{Exactly what accepted 032 adds when $X$ is abelian.}
For an abelian $2m$-fold $X$, 032, Corollary 8.4, asserts that
\[
 (-1)^m\deg_X(z^2)>0
 \quad\text{for every nonzero primitive rational algebraic
 cohomology class }z.
\]
Consequently a primitive cycle which is numerically trivial has zero
cohomology class: its self-intersection is zero, and the displayed
definiteness rules out a nonzero class. Thus 032 explicitly removes
the middle numerical/homological zero-space gate on the primitive
domain of H9-T11. The strict conclusion here is an application of
the accepted theorem; it is not being inferred from semipositivity
on a quotient alone.

One can also derive the whole equality of numerical and homological
equivalence on an abelian variety from exactly the inputs already
accepted. Here are the linear-algebra details. Algebraicity of the
Lefschetz operators makes the primitive decomposition of an algebraic
class algebraic. For a class $a\in\operatorname{Alg}^r(X)$ with
$r\le n/2$, write
\[
 a=\sum_j h^j a_j,
 \qquad a_j\in\operatorname{Alg}^{r-j}(X)_{\rm prim}.
\]
Set the algebraic complementary class
\[
 b=\sum_j(-1)^{r-j}h^{n-2r+j}a_j.
\]
Different primitive strings are orthogonal: if $j<k$, the exponent
$n-2r+j+k$ in the mixed intersection exceeds
$n-2(r-j)$, so it kills $a_j$ by primitivity; interchange the factors
when $k<j$. Hence
\[
 \int_X ab
   =\sum_j(-1)^{r-j}\int_Xa_j^2h^{n-2(r-j)}>0
       \quad\text{if }a\ne0,                               \tag{H9.38}
\]
by 032. A numerically trivial class cannot have this nonzero algebraic
test. For $r>n/2$, put $s=n-r$ and write
$a=h^{n-2s}\sum_j h^j a_j$ using the algebraic inverse of hard
Lefschetz, with $a_j$ primitive of codimension $s-j$. The complementary
algebraic class $b=\sum_j(-1)^{s-j}h^j a_j$ gives
$\int_Xab=\sum_j(-1)^{s-j}\int_Xa_j^2h^{n-2s+2j}>0$
when $a\ne0$, by the same orthogonality and 032. The reverse
implication, homological triviality implies numerical triviality,
is immediate from compatibility of cycle classes and intersection.
This proves the needed equality for abelian varieties rather than
silently inserting it as another hypothesis.

For a general $X$ or for a fixed complete-intersection axis in an
abelian variety, 032 has not been assumed for that variety. Neither
an embedding into an abelian variety nor an unspecified relation of
motives replaces (H9.38). The separate numerical/homological gate
therefore remains unless an actual transfer preserving the relevant
forms is proved.

\paragraph{Relation to the archive and to the height statements.}
Cote 44, scan 9, writes the self-incidence divisor of a pointed family
and asks for positivity on the abelian quotient by its kernel. The
family $z_Z$ is an explicit instance of that construction, with the
primitive qualification recorded in the margin of scan 8. Scan 8
also expressly makes the Hodge reduction dependent on further
unproved assertions. The present argument identifies a precise
conditional numerical reduction; it does not identify those
unwritten assertions historically.

Theorem 5 in scan 4 starts with \emph{algebraically} trivial generic
restrictions, whereas literal H9 starts with homologically trivial
restrictions. H9-T9 proves the height/trace consequence on the former
domain. H9-T11 uses an elliptic family which is algebraically trivial
by construction, so it never crosses that invalid domain step.
It uses polarization directly, not a N\'eron--Tate height on a trace
quotient. In particular, when $X\times E$ is taken constant over a
function field, its constant elliptic family is killed by the trace
quotient and cannot be recovered from that quotient's height.

No assertion here concerns the whole complementary arithmetic Chow
summand $\mathrm{CH}(N)$. Positivity of the separately constructed
Fermat class and its rational span, and positivity of a geometric
incidence quotient, remain different statements with their stated
domains.

\paragraph{Status.}
This is a conditional deduction written in this run, with its complete
intersection and elliptic Poincar\'e calculations included. It is not
claimed to reproduce the unpublished 1964 argument, and it is not an
unconditional proof of the Hodge standard conjecture for arbitrary
varieties. It concerns the incidence polarization itself; it does not
identify a function-field N\'eron--Tate height on a trace quotient with
every primitive middle cycle.
